\section{A polynomial certificate for two zero-free constants}
Let $U_j(\cos t)\sin t=\sin((j+1)t)$ and $T_j(\cos t)=\cos(jt)$, interpreted by their polynomial recurrences at the endpoints. For integer numerators $c_{\ell j}$ put
\[
        A(x,y)=10^{-6}\sum_{\ell,j}c_{\ell j}U_j(x)T_\ell(y).
\]
The normalized target coefficient is $c_{11}/10^6=1$. For the polynomial optimized for the Weyl and Nelson estimates the nonzero integer coefficients are
\[
\begin{array}{c|rrrrrr}
 \ell\backslash j&0&1&2&3&4&6\\\hline
 0&469198&0&417266&0&33789&1678\\
 1&0&1000000&0&255287&0&0\\
 2&459526&0&339719&0&0&908\\
 3&0&194304&0&2181&0&0\\
 4&38741&0&0&0&0&0\\
 6&1135&0&477&0&67&0
\end{array}
\]
\paragraph{Proposition.}
This polynomial is nonnegative on $\mathbb R\times[-1,1]$. Give each term with $\ell j=0$ zero ordinate cost, each moving term with $j=1$ cost $1/3$, and each remaining moving term cost
\[
 d_j=\frac{k}{4}\left(1-\frac2{3k^5-2k^4-k^2}\right),\qquad k=j+1 .
\]
Then its normalized pole weight and total moving cost are
\[
 P=469198/10^6,\qquad G=33956950057987/37263240000000 .
\]
In particular $G/(1-\sqrt P)^2<9.18275$.

\paragraph{Proof.}
For $-1\le x,y\le1$ set $x=(a-b)/(a+b)$ and $y=(u-v)/(u+v)$ and multiply $10^6A$ by $(a+b)^6(u+v)^6$. The resulting homogeneous polynomial has a $7$-by-$7$ integer coefficient array in the monomials $a^i b^{6-i}u^jv^{6-j}$. A nonnegative array proves the desired inequality on the positive orthant. To test a mixed array, partition $u\le v$ and $v\le u$. Substitute $(u,v)=(s,s+2t)$ and $(u,v)=(2s+t,t)$ in the respective pieces. These substitutions preserve nonnegativity and cover the orthant. Alternate the two variable pairs at successive bisections. For the array in the statement the numbers of accepted rectangles at depths $3,4,6,7,8,9,10,11,12,13,14$ are respectively $2,2,12,14,26,44,68,64,112,94,68$. The exterior calculation uses $x=(a+b)/b$ and clearing factor $b^6(u+v)^6$. It terminates on five rectangles at depths at most four. The parity condition $j\equiv\ell\pmod2$ gives $A(-x,-y)=A(x,y)$ for the other exterior component. Lean checks the original polynomial identities, both integer arrays, every bisection and the passage back to real evaluation. The cost is a finite rational sum. An upper rational enclosure of $\sqrt P$ proves the stated strict inequality. See \texttt{EllipticWeyl.lean}, theorems \texttt{fourier\_nonnegative} and \texttt{width\_coefficient}.

\paragraph{Earlier elliptic polynomial.}
For the degree-six table of the preceding manuscript subtract two from its unnormalized static coefficient, replacing $462244$ by $462242$. The same proof certifies this smaller polynomial on $\mathbb R\times[-1,1]$. For its original Stechkin cost the reciprocal width coefficient is now strictly below $12.32544$. At a special place replace $U_j(x)$ by $b^j$ for any signed $b\in[-1,1]$. Subtracting a further $50000$ in numerator units, the homogeneous array is nonnegative after at most four bisections, on ten terminal rectangles. Thus the adjusted local sum is at least $1/20$. A ramified quadratic twist averages the contributions at $b$ and $-b$, preserving this margin. The theorems \texttt{table\_width\_coefficient}, \texttt{special\_margin} and \texttt{special\_even\_margin} give the statements with the exact coefficient lists.

For $1<\sigma\le1.02$ and $\sigma_1\ge\sigma$, the local margin allows the Stechkin explicit formula to retain at least
\[
 \frac{1-\kappa}{20}\sum_{v\ {\rm potentially\ multiplicative}}
           \frac{\log q_v}{q_v^{1.02}-1}.
\]
It is subtracted from the cost side of the enclosure. Indeed $q_v^{-\nu\sigma}-\kappa q_v^{-\nu\sigma_1}\ge(1-\kappa)q_v^{-\nu\sigma}$, and its sum is a geometric series. Identifying the numerical pole and ordinate charges with the Hadamard explicit formula for an actual automorphic function is an analytic input.
